\section{2023 Applied \& Computational Mathematics}

\begin{exercise}
\label{applied:2023:1}
Consider the forward and the centered finite difference formulas

\[
D _ {h} ^ {+} f \left(x _ {0}\right) = \frac {f \left(x _ {0} + h\right) - f \left(x _ {0}\right)}{h}, \tag {1}
\]

\[
D _ {h} ^ {0} f \left(x _ {0}\right) = \frac {f \left(x _ {0} + h\right) - f \left(x _ {0} - h\right)}{2 h}, \tag {2}
\]

to approximate the derivative of $f$ at a point $x _ { 0 }$ . Assume $f$ is a smooth function in a neighborhood of $x _ { 0 }$ containing the points $x _ { 0 } + h$ and $x _ { 0 } - h$ .

(a) Prove that $D _ { h } ^ { + } f ( x _ { 0 } )$ and $D _ { h } ^ { 0 } f ( x _ { 0 } )$ approximate $f ^ { \prime } ( x _ { 0 } )$ to $O ( h )$ and $O ( h ^ { 2 } )$ , respectively.   
(b) Derive an $O ( h ^ { 2 } )$ approximation to $f ^ { \prime } ( x _ { 0 } )$ from $D _ { h } ^ { + } f ( x _ { 0 } )$ by doing Richardson extrapolation.   
(c) Take $f ( x ) = \sin x$ and $x _ { 0 } = 0$ . Prove that both $D _ { h } ^ { + } f ( x _ { 0 } )$ and $D _ { h } ^ { 0 } f ( x _ { 0 } )$ converge quadratically to $f ^ { \prime } ( x _ { 0 } )$ as $h  0$ and that in fact they produce the same approximation to $f ^ { \prime } ( x _ { 0 } )$ in this particular case.
\end{exercise}

\begin{solution}
\textbf{(a) Error Analysis:}
By Taylor's expansion near $x_0$:
\[ f(x_0 + h) = f(x_0) + h f'(x_0) + \frac{h^2}{2} f''(x_0) + \frac{h^3}{6} f'''(x_0) + O(h^4) \]
\[ f(x_0 - h) = f(x_0) - h f'(x_0) + \frac{h^2}{2} f''(x_0) - \frac{h^3}{6} f'''(x_0) + O(h^4) \]
For the forward difference:
\[ D_h^+ f(x_0) = \frac{h f'(x_0) + \frac{h^2}{2} f''(x_0) + O(h^3)}{h} = f'(x_0) + \frac{h}{2} f''(x_0) + O(h^2) = f'(x_0) + O(h). \]
For the centered difference:
\[ D_h^0 f(x_0) = \frac{2 h f'(x_0) + \frac{h^3}{3} f'''(x_0) + O(h^5)}{2h} = f'(x_0) + \frac{h^2}{6} f'''(x_0) + O(h^4) = f'(x_0) + O(h^2). \]

\textbf{(b) Richardson Extrapolation:}
Let $D_h^+ = f'(x_0) + c_1 h + c_2 h^2 + \dots$. Then $D_{2h}^+ = f'(x_0) + 2 c_1 h + 4 c_2 h^2 + \dots$.
To eliminate the $O(h)$ term, compute $2 D_h^+ - D_{2h}^+$:
\[ R(h) = 2 D_h^+ f(x_0) - D_{2h}^+ f(x_0) = (2 f' + 2c_1 h + 2c_2 h^2) - (f' + 2c_1 h + 4c_2 h^2) + O(h^3) \]
\[ R(h) = f' - 2c_2 h^2 + O(h^3) = f'(x_0) + O(h^2). \]
Specifically, $R(h) = \frac{2}{h}(f(x_0+h)-f(x_0)) - \frac{1}{2h}(f(x_0+2h)-f(x_0)) = \frac{-f(x_0+2h) + 4f(x_0+h) - 3f(x_0)}{2h}$.

\textbf{(c) Case $f(x) = \sin x, x_0 = 0$:}
Here $f'(0) = \cos 0 = 1$ and $f''(0) = -\sin 0 = 0$.
Forward difference: $D_h^+ \sin(0) = \frac{\sin h - 0}{h} = \frac{\sin h}{h}$.
Centered difference: $D_h^0 \sin(0) = \frac{\sin h - \sin(-h)}{2h} = \frac{2\sin h}{2h} = \frac{\sin h}{h}$.
Since $f''(0) = 0$, the $O(h)$ term in the error of $D_h^+$ vanishes, and both yield $\frac{\sin h}{h} = 1 - \frac{h^2}{6} + O(h^4)$, which is $O(h^2)$ accurate.
\end{solution}

\begin{exercise}
\label{applied:2023:2}
For functions defined on a closed interval $[ 0 , 1 ]$ , we want to compute the following definite integral,

\[
I [ f ] = \int_ {0} ^ {1} f (x) \log (1 / x) d x.
\]

Here we consider the weight function $\log ( 1 / x )$ , and denote $P _ { n } ( x )$ as the monic orthogonal polynomials for the corresponding weighted inner product.

(a) Let $P _ { 0 } = 1$ . Find $P _ { 1 } ( x )$ , and the corresponding node $x _ { 1 } ^ { 1 }$ and weight $\omega _ { 1 } ^ { 1 }$ for the 1-point Gaussian quadrature rule.   
(b) Derive a recursive formula for $P _ { n + 1 } ( x )$ using $P _ { n } ( x )$ and $P _ { n - 1 } ( x )$ .   
(c) Consider the normalized orthogonal polynomials $Q _ { n } ( x ) = P _ { n } ( x ) / \vert \vert P _ { n } \vert \vert$ , where

\[
| | P _ {n} | | = \sqrt {P _ {n} (x) ^ {2} \log (1 / x) d x}.
\]

Derive a recursive formula for $Q _ { n + 1 } ( x )$ using $Q _ { n } ( x )$ and $Q _ { n - 1 } ( x )$ .

(d) Use the above recursive formula to show that $x = \lambda$ is a node of the 4-point Gaussian quadrature if and only if it is an eigenvalue of a symmetric, tridiagonal matrix. Write out the form of the symmetric and tridiagonal matrix explicitly.
\end{exercise}

\begin{solution}
\textbf{(a) Gaussian Quadrature:}
Moments $\mu_k = \int_0^1 x^k \log(1/x) dx$. Let $x = e^{-t}$: $\mu_k = \int_0^\infty e^{-kt} t e^{-t} dt = \int_0^\infty t e^{-(k+1)t} dt = \frac{1}{(k+1)^2}$.
$\mu_0 = 1, \mu_1 = 1/4$.
Monic $P_1(x) = x - c$. Orthogonality $\langle P_1, P_0 \rangle = 0 \implies \mu_1 - c \mu_0 = 0 \implies c = 1/4$.
$P_1(x) = x - 1/4$. Node $x_1^1 = 1/4$.
Weight $\omega_1^1 = \mu_0/P_0(x_1^1) = 1$.

\textbf{(b) Recurrence for Monic Polynomials:}
Monic orthogonal polynomials satisfy $P_{n+1}(x) = (x - \alpha_n) P_n(x) - \beta_n P_{n-1}(x)$, where:
\[ \alpha_n = \frac{\langle x P_n, P_n \rangle}{\langle P_n, P_n \rangle}, \quad \beta_n = \frac{\langle P_n, P_n \rangle}{\langle P_{n-1}, P_{n-1} \rangle}. \]

\textbf{(c) Recurrence for Normalized Polynomials:}
Let $b_n = \sqrt{\beta_n} = \frac{\|P_n\|}{\|P_{n-1}\|}$. Substituting $P_n = \|P_n\| Q_n$ into (b):
\[ \|P_{n+1}\| Q_{n+1} = (x - \alpha_n) \|P_n\| Q_n - \|P_n\|^2/\|P_{n-1}\|^2 \|P_{n-1}\| Q_{n-1} \]
Dividing by $\|P_{n+1}\|$:
\[ Q_{n+1} = \frac{\|P_n\|}{\|P_{n+1}\|} (x - \alpha_n) Q_n - \frac{\|P_n\|}{\|P_{n+1}\|} b_n Q_{n-1} \]
\[ b_{n+1} Q_{n+1} = (x - \alpha_n) Q_n - b_n Q_{n-1} \implies x Q_n = b_{n+1} Q_{n+1} + \alpha_n Q_n + b_n Q_{n-1}. \]

\textbf{(d) Jacobi Matrix:}
The nodes are the roots of $P_4(x)$ (and thus $Q_4(x)$). The recurrence for $n=0,1,2,3$ with $Q_{-1}=0$ in vector form is:
\[ x \begin{pmatrix} Q_0 \\ Q_1 \\ Q_2 \\ Q_3 \end{pmatrix} = \begin{pmatrix} \alpha_0 & b_1 & 0 & 0 \\ b_1 & \alpha_1 & b_2 & 0 \\ 0 & b_2 & \alpha_2 & b_3 \\ 0 & 0 & b_3 & \alpha_3 \end{pmatrix} \begin{pmatrix} Q_0 \\ Q_1 \\ Q_2 \\ Q_3 \end{pmatrix} + \begin{pmatrix} 0 \\ 0 \\ 0 \\ b_4 Q_4(x) \end{pmatrix} \]
At a root $x = \lambda_j$ of $Q_4$, $\lambda_j$ is an eigenvalue of the symmetric tridiagonal Jacobi matrix $J_4$.
\end{solution}

\begin{exercise}
\label{applied:2023:3}
Let $A$ be a real $n \times n$ matrix with distinct eigenvalues such that

\[
\left| \lambda_ {1} \right| > \left| \lambda_ {2} \right| \geq \left| \lambda_ {3} \right| \geq \dots \geq \left| \lambda_ {n} \right| \geq 0,
\]

with corresponding eigenvectors $\{ v _ { j } \} _ { j = 1 } ^ { n }$ .

(a) Show that the power iteration

\[
z _ {m} = \frac {A ^ {m} z _ {0}}{| | A ^ {m} z _ {0} | | _ {\infty}} \longrightarrow \pm \frac {v _ {1}}{| | v _ {1} | | _ {\infty}}, \quad \forall z _ {0} \in \mathbb {R} ^ {n}.
\]

(b) Consider the following iteration with initial guess $x _ { 0 } = y _ { 0 } = 1$

\[
x _ {n + 1} = x _ {n} + y _ {n}, \quad y _ {n + 1} = x _ {n + 1} + x _ {n}.
\]

Show that $y _ { n } / x _ { n } \to { \sqrt { 2 } }$ as $n \to \infty$
\end{exercise}

\begin{solution}
\textbf{(a) Power Method:}
Let $z_0 = \sum_{j=1}^n c_j v_j$. $A^m z_0 = \lambda_1^m [c_1 v_1 + \sum_{j=2}^n c_j (\lambda_j/\lambda_1)^m v_j]$.
Since $|\lambda_j/\lambda_1| < 1$, $(\lambda_j/\lambda_1)^m \to 0$. Thus $A^m z_0 \approx c_1 \lambda_1^m v_1$.
Normalizing gives $z_m \to \frac{c_1 \lambda_1^m v_1}{\|c_1 \lambda_1^m v_1\|_\infty} = \text{sgn}(c_1 \lambda_1^m) \frac{v_1}{\|v_1\|_\infty}$.

\textbf{(b) Ratio Convergence:}
The system is $\begin{pmatrix} x_{n+1} \\ y_{n+1} \end{pmatrix} = \begin{pmatrix} 1 & 1 \\ 2 & 1 \end{pmatrix} \begin{pmatrix} x_n \\ y_n \end{pmatrix}$.
Eigenvalues of $M = \begin{pmatrix} 1 & 1 \\ 2 & 1 \end{pmatrix}$ are $1 \pm \sqrt{2}$.
Dominant eigenvalue $\lambda_1 = 1+\sqrt{2}$ has eigenvector $v_1 = \begin{pmatrix} 1 \\ \sqrt{2} \end{pmatrix}$.
By the power method, the state vector aligns with $v_1$, so $y_n/x_n \to \sqrt{2}$.
\end{solution}

\begin{exercise}
\label{applied:2023:4}
Consider the initial value problem

\[
y ^ {\prime} = f (t, y), \quad 0 <   t \leq T. \tag {3}
\]

\[
y (0) = y _ {0}. \tag {4}
\]

Assume $f$ is continuous and Lipschitz in $y$ in $[ 0 , T ] \times ( - \infty , \infty )$ . Denote $y _ { n } \approx y ( t _ { n } )$ , $t _ { n } = n h$ , and $h = T / N$ , with $N$ a positive integer, and consider the one-step method

\[
y _ {n + 1} = y _ {n} + \alpha h f (t _ {n}, y _ {n}) + \beta h f (t _ {n} + \gamma h, y _ {n} + \gamma h f (t _ {n}, y _ {n})),
\]

where $\alpha$ , $\beta$ and $\gamma$ are real parameters.

(a) Prove that the method is consistent if and only if $\alpha + \beta = 1$ , and the order of the method can not exceed 2.   
(b) Suppose that a second-order method of the above form is applied to $f ( t , y ) = - \lambda y$ with $\lambda > 0$ , and the initial condition $y _ { 0 } = 1$ . Show that the sequence $( y _ { n } ) _ { n \geq 0 }$ is bounded if and only if $\begin{array} { r } { h \leq \frac { 2 } { \lambda } } \end{array}$ . Show further that for such $h$ ,

\[
| y (t _ {n}) - y _ {n} | \leq \frac {1}{6} \lambda^ {3} h ^ {2} t _ {n}, \quad n \geq 0.
\]
\end{exercise}

\begin{solution}
\textbf{(a) Consistency and Order:}
LHS: $y(t+h) - y(t) = h f + \frac{h^2}{2} (f_t + f f_y) + O(h^3)$.
RHS: $(\alpha+\beta) h f + \beta \gamma h^2 (f_t + f f_y) + O(h^3)$.
Consistency: $h f = (\alpha+\beta) h f \implies \alpha+\beta=1$.
Order 2: $\beta \gamma = 1/2$. Order 3 would requires matching terms like $f_y^2 f$, but this is a 2-stage RK method. The maximum order of an $s$-stage RK method is $s$, so $s=2$.

\textbf{(b) Linear Stability:}
$y_{n+1} = y_n [1 - h\lambda + \frac{1}{2} (h\lambda)^2] \equiv g(h\lambda) y_n$.
Bounded iff $|1 - z + z^2/2| \leq 1$ for $z=h\lambda$. This yields $z \in [0, 2]$, so $h \leq 2/\lambda$.
Error $e_n = |y(t_n) - y_n|$. $e_{n+1} \leq g(z) e_n + |y(t_{n+1}) - g(z)y(t_n)|$.
$LTE \approx |e^{-z} - (1-z+z^2/2)| \leq \frac{z^3}{6}$ (Taylor remainder).
Then $e_n \leq \sum_{k=0}^{n-1} \frac{(h\lambda)^3}{6} = n \frac{1}{6} \lambda^3 h^3 = \frac{1}{6} \lambda^3 h^2 t_n$.
\end{solution}

\begin{exercise}
\label{applied:2023:5}
Let $u ( t , x )$ be the solution of the initial-boundary value problem

\[
u _ {t} = D u _ {x x}, \quad 0 <   x <   L, \quad 0 <   t \leq T, \tag {5}
\]

\[
u (0, x) = f (x) \tag {6}
\]

\[
u (t, 0) = u (t, L) = 0, \tag {7}
\]

where $L > 0$ and $D > 0$ . Consider the finite difference scheme

\[
\frac {u _ {j} ^ {n + 1} - u _ {j} ^ {n}}{\Delta t} = D \frac {u _ {j + 1} ^ {n} - 2 u _ {j} ^ {n} + u _ {j - 1} ^ {n}}{(\Delta x) ^ {2}}, \quad j = 1, \dots , M - 1, \quad n = 0, 1, \dots , N - 1 \tag {8}
\]

with $u _ { 0 } ^ { n } = u _ { M } ^ { n } = 0$ for all $n$ and $u _ { j } ^ { \mathrm { 0 } } = f ( j \Delta x )$ , $j = 0 , \ldots , M$ . Here $\Delta t = T / N$ and $\Delta x = L / M$ and $u _ { j } ^ { n } \approx u ( n \Delta t , j \Delta x )$ .

(a) Prove that (8) is consistent with (5).   
(b) Prove that if $\begin{array} { r } { \Delta t { } \leq \frac { 1 } { 2 D } ( \Delta x ) ^ { 2 } } \end{array}$ the finite difference scheme (8) is stable under the $l ^ { \infty }$ norm.   
(c) Prove that if $\begin{array} { r } { \Delta t { } \leq \frac { 1 } { 2 D } ( \Delta x ) ^ { 2 } } \end{array}$ the finite difference scheme (8) converges in the $l ^ { \infty }$ norm to the exact solution of (5)-(7).
\end{exercise}

\begin{solution}
\textbf{(a) Consistency:}
Taylor expansion of the discrete operator applied to the exact solution $U$:
\[ \frac{U_j^{n+1}-U_j^n}{\Delta t} - D \frac{U_{j+1}^n - 2U_j^n + U_{j-1}^n}{(\Delta x)^2} = [U_t + \frac{\Delta t}{2} U_{tt}] - D [U_{xx} + \frac{\Delta x^2}{12} U_{xxxx}] + \dots \]
Since $U_t = D U_{xx}$, the error is $O(\Delta t + \Delta x^2)$, which vanishes as $\Delta t, \Delta x \to 0$.

\textbf{(b) Stability:}
$u_j^{n+1} = (1-2r) u_j^n + r u_{j+1}^n + r u_{j-1}^n$ with $r = \frac{D \Delta t}{\Delta x^2}$.
If $r \leq 1/2$, then $1-2r \geq 0$ and all coefficients are non-negative.
$|u_j^{n+1}| \leq (1-2r+r+r) \|u^n\|_\infty = \|u^n\|_\infty$.

\textbf{(c) Convergence:}
The error $e^n$ satisfies $e^{n+1} = M e^n + \Delta t \tau^n$ where $\|M\|_\infty = 1$ and $\|\tau^n\|_\infty = O(\Delta t + \Delta x^2)$.
By induction, $\|e^n\|_\infty \leq n \Delta t \max \|\tau\| = t_n O(\Delta t + \Delta x^2)$.
\end{solution}

\begin{exercise}
\label{applied:2023:6}
Let $\psi ^ { \varepsilon } ( t , x )$ be the solution to the following Schr¨odinger equation:

\[
\mathrm {i} \varepsilon \frac {\partial \psi^ {\varepsilon}}{\partial t} = - \frac {\varepsilon^ {2}}{2} \nabla_ {x} ^ {2} \psi^ {\varepsilon} + V (x) \psi^ {\varepsilon}, x = (x _ {1}, \dots , x _ {n}) ^ {\mathrm {T}} \in \mathbb {R} ^ {n},
\]

where $\mathrm { i } = \sqrt { - 1 }$ , $\varepsilon \ll 1$ is a small positive real number (rescaled Planck constant), $\nabla _ { x } ^ { 2 } = \sum _ { j = 1 } ^ { n } \partial _ { x _ { j } } ^ { 2 }$ n ∂2xj , and $V ( x ) \in C ^ { \infty } ( \mathbb { R } ^ { n } )$ is the potential function.

Consider the WKB expansion

\[
\psi^ {\varepsilon} (t, x) = A (t, x) e ^ {\mathrm {i} \frac {S (t , x)}{\varepsilon}},
\]

(a) Derive equations for $A ( t , x )$ and $S ( t , x )$ by asymptotic expansion. (Here both $A ( t , x )$ and $S ( t , x )$ are real-valued functions, and do not depend on $\varepsilon$ .)   
(b) Define $u ( t , x ) = \nabla _ { x } S ( t , x ) \in \mathbb { R } ^ { n }$ . Derive an equation for $u ( t , x )$ . Suppose $u ( 0 , x ) \in$ $C ^ { \infty } ( \mathbb { R } ^ { n } )$ , will $u ( t , x )$ always be in $C ^ { \infty } ( \mathbb { R } ^ { n } )$ for all $t > 0$ ? Explain why.
\end{exercise}

\begin{solution}
\textbf{(a) WKB Equations:}
Substituting $\psi = A e^{iS/\epsilon}$ and equating powers of $\epsilon$:
$O(1)$: $-S_t = \frac{1}{2} |\nabla S|^2 + V$, which is the Hamilton-Jacobi equation.
$O(\epsilon)$: $A_t + \nabla S \cdot \nabla A + \frac{1}{2} A \Delta S = 0$, which is the transport equation.

\textbf{(b) Gradient Equation:}
Let $u = \nabla S$. $\nabla(S_t + \frac{1}{2} |\nabla S|^2 + V) = u_t + (u \cdot \nabla)u + \nabla V = 0$.
This is a pressureless Euler-type equation. It does not stay in $C^\infty$ as trajectories (characteristics) can cross, leading to shock formation (gradient catastrophe) in finite time.
\end{solution}
